% !TEX root = ../main.tex

\chapter{System Design and Technical Mechanisms}

\section{Design Method and Implementation Cutoff}

This chapter translates the four technical challenges introduced in Chapter~1 into mechanisms that can be inspected and evaluated. The design follows one dependency chain rather than four isolated feature lists. Revision-safe representation (C1) supplies the identifiers and provenance required by evidence-bounded assistance (C2). The same revision-scoped representations provide inputs to the interest model (C3), while the mobile integration and citation graph (C4) expose evidence, feedback, and failure state to the reader. Each section therefore states a problem, the implemented mechanism, a considered alternative, and the evaluation hook that Chapter~4 must exercise.

Unless a subsection is explicitly marked \planned{}, implementation claims in this chapter refer to the integrated repository snapshot \texttt{b945d3c}. This cutoff contains the backend pipeline, behavior baseline, citation-graph platform, Android graph flow, and durable behavioral-event upload path. A passing test is treated as evidence that the checked-out software satisfies a specified invariant; it is not treated as a measurement of research quality, user benefit, latency, or robustness in an uncontrolled environment. Those outcomes remain \pending{} until the corresponding E1--E4 protocols are executed.

MNEME has four operational trust regions. The Android client owns presentation, cached reader state, background synchronization, and local interaction capture. The FastAPI application owns authentication, request validation, domain services, and public state transitions. PostgreSQL is authoritative for papers, revisions, artifacts, jobs, preferences, events, digests, and resolved citation relations; Redis and ARQ execute attempts but do not define durable truth. External scholarly and model providers remain outside the trusted boundary. Identifiers, configuration versions, quality state, and failures must remain explicit whenever data cross one of these regions.

\begin{figure}[htbp]
\centering
\fbox{\parbox{0.88\textwidth}{
\textbf{FIGURE PLACEHOLDER --- Evidence-first mobile literature loop.}\\
The final vector figure will show: scholarly source $\rightarrow$ exact-revision ingestion $\rightarrow$ versioned artifacts and chunks $\rightarrow$ grounded answer and briefing $\rightarrow$ mobile inspection and behavioral events $\rightarrow$ versioned interest update. Trust boundaries will separate Android, FastAPI workers, PostgreSQL/Redis, and external providers.\\
\textbf{Owner:} Hanyang (visual integration), with Ruiyu and Yifan validating technical labels. \textbf{Dependency:} final API and worker state diagrams. \textbf{Expected artifact:} editable vector figure plus print-safe PDF. \textbf{Finalization trigger:} implementation freeze for E1--E4.
}}
\caption{Evidence-first dependency loop and trust boundaries. The final figure will distinguish implemented paths from planned final-version extensions.}
\label{fig:loop}
\end{figure}

Figure~\ref{fig:loop} will serve as the chapter's visual index. Three invariants connect its components. First, a paper-derived output must identify the exact source revision. Second, a state-changing request must carry a stable identity so that retry does not duplicate its effect. Third, degraded results and failures must remain visible instead of being converted into apparently complete output. These invariants turn the central problem of continuous literature understanding into testable system properties.

\section{C1: Revision-Safe Literature Representation}
\label{sec:c1-revision}

\subsection{Logical Works, Observed Revisions, and Artifact Lineage}

The first technical challenge is continuity under change: a paper identifier remains useful across updates, but an answer or embedding is valid only for the content from which it was derived. MNEME therefore separates a logical paper identifier \(p\) from an observed revision identifier \(r\). A \texttt{papers} row represents the work and maintains discovery metadata; a \texttt{paper\_versions} row records an exact provider revision, content checksum, storage location, parser quality, and processing state. Downstream chunks, embeddings, summaries, jobs, and graph synchronization inputs bind to the revision row rather than only to the logical work.

The lineage of a derived artifact \(a\) is modeled as
\[
  a = f(p,r,h,\theta),
\]
where \(h\) is the source-content checksum and \(\theta\) includes the stage, parser, prompt, model, or embedding version relevant to transformation \(f\). Reuse is safe only when all material inputs for that artifact class match. In particular, equality of \(p\) does not imply equality of \(a\). The database constraints and repositories enforce this distinction by using internal UUIDs for relationships and uniqueness constraints for provider identifiers and revision labels.

The arXiv synchronizer applies this identity model in two steps. It normalizes a provider identifier to locate or create the logical paper, then inserts only previously unseen revisions. A newer metadata observation may refresh the logical paper snapshot without mutating an existing revision. The transaction also normalizes author identities and returns exact revision identifiers for subsequent work. This design avoids an in-place-update alternative in which the latest PDF silently replaces its predecessor; that alternative is simpler to query but makes prior summaries and citations impossible to interpret after an update.

\subsection{Durable Stage Identity and Recovery}

Processing is decomposed into metadata fetch, PDF download, parsing, summarization, chunking, embedding, and digest assembly. Paper-scoped stages carry both \texttt{paper\_id} and \texttt{paper\_version\_id}. A canonical serialization of the stage, scope, pipeline version, and material inputs is hashed into an idempotency key. A unique database constraint then maps logically equivalent submissions to one durable job row even when an API request, scheduler, or worker retries.

PostgreSQL records the job state \texttt{queued}, \texttt{running}, \texttt{succeeded}, or \texttt{failed}, along with attempt counts, timestamps, leases, and the latest internal error. Redis/ARQ receives a stable attempt identifier only after the durable row exists. Dispatch leases allow interrupted enqueue operations to be reclaimed, and a recovery service periodically re-enqueues revision-scoped work that is eligible for retry. The public \texttt{GET /jobs/\{job\_id\}} endpoint exposes durable progress while withholding internal exception text. This database-first design was chosen over treating the queue as the source of truth because queue retention and worker availability should not determine whether processing history can be reconstructed.

The download and artifact stores extend idempotency to the file system. Downloads enforce size, media signature, retry policy, and an expected checksum when available. New bytes are written to a temporary file, synchronized, and atomically installed with \texttt{os.replace}. Structured parse artifacts and their sidecars follow the same installation pattern. Consequently, a process failure can leave either the previous complete artifact or a recoverable temporary file, rather than a partially overwritten file that appears valid.

\subsection{Parsing, Quality Tiers, and Degraded Operation}

The parser combines PyMuPDF text blocks with pdfplumber layout hints. It removes repeated headers and footers, estimates column structure, orders blocks within each column, and detects headings before assembling sections. This approach retains more source structure than flat text extraction while remaining deterministic and inspectable. Scientific PDFs nevertheless vary in embedded text, reading order, typography, and page layout. The parser therefore records one of three implemented quality tiers:
\begin{itemize}
  \item \texttt{structured}: section structure and ordered text pass the parser's checks;
  \item \texttt{text\_only}: useful full text is available but structural confidence is insufficient;
  \item \texttt{abstract\_only}: full-text processing cannot be trusted, so downstream services use provider metadata rather than pretending that section-level evidence exists.
\end{itemize}
The processing status is reconciled from artifacts belonging to the latest observed revision. A revision can become \texttt{ready} only when a non-abstract summary exists, chunks exist, and those chunks have compatible embeddings; otherwise the paper remains partial or failed. Quality is therefore a control input to downstream behavior, not a decorative label.

\subsection{Revision Isolation and the E1 Contract}

Section-aware chunking runs only after the parser has produced a revision-scoped document. Each chunk stores the exact \texttt{paper\_version\_id}, section title, stable local order, content hash, token count, and embedding provenance. Sentence-aware splitting and bounded overlap avoid crossing section boundaries when the source structure permits it. Chunking research indicates that retrieval performance depends on the interaction between segmentation, context expansion, and later retrieval stages rather than on a universally optimal chunk size~\cite{wang2024bestpractices}. MNEME therefore treats its current section-aware policy as an implemented configuration to be compared in E2, not as an already proven improvement.

E1 closes the C1 argument by testing the properties that implementation tests alone cannot establish across a representative corpus. Its fixtures will include multiple revisions, malformed or layout-diverse PDFs, interrupted stages, duplicate submissions, and stale-artifact scenarios. The primary outcomes are cross-revision contamination, idempotent recovery, parse-quality distribution, provenance completeness, and processing latency. The current repository already tests key software invariants, including revision-scoped retrieval and job idempotency; corpus-level rates and performance values remain \pending{}. This bounded conclusion is important: C1 establishes an inspectable mechanism for continuity, while E1 must determine how reliably that mechanism operates on the selected evidence set.

\section{C2: Evidence-Bounded Assistance}

\subsection{Retrieval Scope and Trace}

Each question has a declared scope: one revision, a selected set of papers, or a user library snapshot.
The retrieval trace stores the query, embedding model, candidate filters, ranking configuration, retrieved chunk identifiers, scores, and context expansion.
This trace is necessary because the same generated answer may be interpreted differently when produced from a single paper versus a broad library.

Long context is not treated as an automatic substitute for retrieval.
Evidence position can affect model use~\cite{liu2024lost}, and retrieval may add irrelevant context.
The planned pipeline therefore evaluates both recall of relevant evidence and resistance to distractors.

\subsection{Claim and Evidence Objects}

The generator returns a structured answer consisting of claims:
\[
  A = \{c_1,c_2,\ldots,c_n\}.
\]
Each claim \(c_i\) contains claim text, one or more evidence references, an answerability state, and verification results.
An evidence reference contains the logical paper and revision identifiers, section and chunk identifiers, a character or sentence span, and the retrieved text hash.

This representation separates four questions:
\begin{enumerate}
  \item \textbf{Identifier validity}: does the cited paper revision exist in the retrieval trace?
  \item \textbf{Evidence availability}: can the cited source span be resolved?
  \item \textbf{Consistency}: does the evidence support, contradict, or fail to determine the claim?
  \item \textbf{Coverage}: do the cited spans support all material parts of the answer?
\end{enumerate}
ALCE motivates explicit citation-quality dimensions~\cite{gao2023alce}, while fine-grained attribution work shows why a document identifier alone is insufficient~\cite{huang2024grounded}.
RAGTruth further demonstrates that RAG outputs can contain unsupported or contradictory content despite retrieved context~\cite{niu2024ragtruth}.

\subsection{Generation and Verification}

The planned generation prompt provides retrieved chunks with immutable metadata and requires structured output.
Post-generation checks reject identifiers that were not retrieved and unresolved spans.
Consistency and coverage checks may combine deterministic validation, natural-language inference, and human review.
No automatic judge is treated as ground truth.
RAGAS and ARES illustrate component-level evaluation and the value of human calibration~\cite{es2024ragas,saadfalcon2024ares}.

If no retrieved evidence supports an answer, the correct behavior is abstention or a scoped clarification request.
The interface should distinguish ``no evidence found'' from provider failure and from a failed consistency check.
This distinction prevents the absence of an answer from being presented as a confident negative claim.

\subsection{Multi-Turn Boundaries}

A follow-up question may depend on prior entities or scope, but conversation history must not silently widen the evidence set.
MNEME retains a bounded number of turns and records the retrieval query derived from each turn.
Scope changes are visible.
E2 treats first-turn, follow-up, non-standalone, and unanswerable questions as separate strata, following the need for dedicated multi-turn RAG evaluation~\cite{katsis2025mtrag}.

\section{C3: Adaptive and Inspectable Interest Memory}

\subsection{Event Semantics}

The planned event vocabulary includes explicit topic edits and paper saves as well as implicit opens, reading duration, skips after exposure, shares, and questions.
Each event contains an event identifier, user identifier, paper and revision identifiers when applicable, event type, timestamp, context, and client-generated deduplication key.
The event model records observation, not intent.
For example, a long open duration may indicate careful reading, interruption, or abandonment.
Its weight must therefore be bounded and tested.

\subsection{Multi-Timescale State}

Let \(e_t\) be the embedding associated with an eligible event and \(w_t\) its declared signal weight.
A baseline short-term state can be represented by
\[
  s_t = \operatorname{norm}\left(\alpha_s s_{t-1} + (1-\alpha_s)w_t e_t\right),
\]
while a long-term state uses a slower update,
\[
  \ell_t = \operatorname{norm}\left(\alpha_\ell \ell_{t-1} + (1-\alpha_\ell)w_t e_t\right),
  \qquad \alpha_\ell > \alpha_s.
\]
The recommendation state may combine explicit preferences \(x_t\), short-term state, and long-term state:
\[
  u_t = \operatorname{norm}\left(\beta_x x_t + \beta_s s_t + \beta_\ell \ell_t\right).
\]
These equations specify a baseline to be frozen before E3.
The coefficients are not claimed optimal.
They must be documented, ablated, and tested against static and recent-only alternatives.

\subsection{Negative and Missing Feedback}

A skip is eligible as negative evidence only after the system establishes exposure.
Silence outside an exposure window is missing data, not a negative preference.
Unpaired-feedback research highlights the difficulty of interpreting engagement and silent disengagement over long horizons~\cite{zhang2025unpaired}.
MNEME therefore stores event provenance and caps the influence of any single implicit event.
Explicit removal of a topic overrides inferred positive evidence until the user changes that decision or an explicit policy expires.

\subsection{Recommendation Score and Reason}

A candidate score can combine semantic match, freshness, followed-author or citation relations, diversity, and delivery constraints:
\[
  S(p,u_t)=\lambda_m M(p,u_t)+\lambda_f F(p)+\lambda_a A(p)
            +\lambda_d D(p,\mathcal{L})-\lambda_n N(p),
\]
where \(M\) is the match to the interest state, \(F\) is freshness, \(A\) is an explicit author or relation boost, \(D\) rewards diversity relative to the current list \(\mathcal{L}\), and \(N\) represents notification or repetition penalties.
Every term must be defined and exposed to E3 ablation.

The displayed recommendation reason is derived from inspectable contributors, such as an explicit topic, a saved paper, or an author watchlist.
It is not a free-form post-hoc explanation of an opaque score.
Editable natural-language profiles demonstrate one route to scrutable preference representations~\cite{ramos2024scrutable}; MNEME adopts the more general requirement that visible reasons and editable state must remain consistent.

\section{C4: Bounded Mobile Integration}

\subsection{Local State and Synchronization}

The planned Android client keeps digest snapshots, paper metadata, reading state, pending events, and user preferences in local storage.
Background work uploads events and checks for new digests subject to operating-system constraints.
Every synchronization operation is idempotent.
Conflicts are resolved by explicit version and timestamp policies rather than last-write-wins without inspection.

\subsection{Notifications}

Notifications are delivery mechanisms, not evidence that a recommendation is useful.
The design supports scheduled digests, bounded high-relevance alerts, deduplication, daily caps, and user-configurable preferences.
The notification contains enough context to explain why it was sent and deep-links to the corresponding digest or paper snapshot.
E4 must measure delivery reliability and task consequences rather than assuming that more timely delivery is beneficial.

\subsection{Graph Boundary}

The planned graph is rendered in a WebView with a Kotlin--JavaScript bridge.
Graph data are supplied as versioned nodes and edges; interaction callbacks identify the selected node and requested action.
The bridge is a trust boundary because injected data, event routing, and gesture conflicts can fail independently.
The UI must display selected state and a stable primary action.
The formative presentation's graph finding motivates these requirements but does not prove the redesign succeeds.

\subsection{Visible Failure and Degraded Operation}

The client distinguishes unavailable content, stale content, provider failure, unsupported answer, and synchronization delay.
Cached content carries its last-updated revision and time.
An offline reader can inspect already downloaded evidence, but the system must not imply that a new question has been verified when the required service is unavailable.

\section{Cross-Cutting Safety and Reproducibility}

The system logs model and prompt versions, retrieval configuration, data revision, and verification policy for each generated output.
Logs exclude unnecessary paper content and behavioral details.
Users can inspect and delete their event history.
Daily budget limits constrain model calls, and content/configuration hashes permit safe caching.
These controls are part of the design; their implementation and effectiveness remain to be verified against a source revision and E1--E4 artifacts.
